% ============================================================================
% lamet-agent — 完整技术文档
% Agentic Workflow for Large-Momentum Effective Theory Analysis
% ============================================================================
\documentclass[12pt,a4paper]{ctexart}

% ---------- Packages ----------
\usepackage[hmargin=2.5cm,vmargin=2.5cm]{geometry}
\usepackage{hyperref}
\usepackage{graphicx}
\usepackage{listings}
\usepackage{xcolor}
\usepackage{enumitem}
\usepackage{booktabs}
\usepackage{tabularx}
\usepackage{caption}
\usepackage{fancyhdr}
\usepackage{titling}
\usepackage{tcolorbox}
\usepackage{amsmath,amssymb}
\usepackage{fontspec}
\setmonofont{DejaVu Sans Mono}[Scale=0.88]

% ---------- Code Listing Style ----------
\definecolor{codebg}{RGB}{245,245,245}
\definecolor{codeframe}{RGB}{200,200,200}
\definecolor{keyword}{RGB}{0,0,180}
\definecolor{string}{RGB}{0,128,0}
\definecolor{comment}{RGB}{128,128,128}

\lstdefinestyle{pycode}{
    backgroundcolor=\color{codebg}, frame=single, rulecolor=\color{codeframe},
    basicstyle=\footnotesize\ttfamily, keywordstyle=\color{keyword}\bfseries,
    stringstyle=\color{string}, commentstyle=\color{comment}\itshape,
    breaklines=true, showstringspaces=false, numbers=left,
    numberstyle=\tiny\color{gray}, numbersep=5pt, tabsize=4, language=Python,
}
\lstdefinestyle{bashcode}{
    backgroundcolor=\color{codebg}, frame=single, rulecolor=\color{codeframe},
    basicstyle=\footnotesize\ttfamily, breaklines=true, showstringspaces=false,
    numbers=left, numberstyle=\tiny\color{gray}, numbersep=5pt, language=bash,
}
\lstdefinestyle{jsoncode}{
    backgroundcolor=\color{codebg}, frame=single, rulecolor=\color{codeframe},
    basicstyle=\footnotesize\ttfamily, breaklines=true, showstringspaces=false,
    numbers=left, numberstyle=\tiny\color{gray}, numbersep=5pt,
}

\tcbuselibrary{skins,breakable}

% ---------- Metadata ----------
\title{{\Huge\bfseries lamet-agent}\\[4pt]
       {\large 完整技术文档}\\[4pt]
       {\normalsize Agentic Workflow for Large-Momentum Effective Theory Analysis}}
\author{}
\date{2026年7月20日}

\pagestyle{fancy}
\fancyhf{}
\fancyhead[L]{\small lamet-agent 技术文档}
\fancyhead[R]{\small \thepage}
\renewcommand{\headrulewidth}{0.4pt}

\begin{document}

\maketitle
\thispagestyle{empty}

\vspace{2cm}
\begin{center}
\begin{tcolorbox}[width=0.85\textwidth,boxrule=0.8pt,colback=blue!3!white,colframe=blue!40!black]
\centering
\textbf{项目地址：}\url{https://github.com/Greyyy-HJC/lamet-agent.git}

\textbf{许可证：}MIT

\textbf{Python 版本要求：}3.10+
\end{tcolorbox}
\end{center}

\newpage
\tableofcontents
\newpage

% ============================================================================
\section{项目概述}
% ============================================================================

\texttt{lamet-agent} 是一个 Python AI 代理框架，专为大动量有效理论（LaMET）格点 QCD 分析设计。它将完整的 LaMET 分析工作流自动化，从关联函数数据出发，经过重整化、傅里叶变换、微扰匹配和外推，最终产生物理分布函数（PDF、DA、TMD）。

\subsection{核心设计理念}

\begin{itemize}[itemsep=4pt]
    \item \textbf{清单驱动（Manifest-Driven）}：通过 JSON 清单定义全局数据源池和分阶段任务列表
    \item \textbf{DAG 任务图}：任务 ID 形成有向无环图，关联函数任务分组原始数据集，后续任务通过角色命名的输入消费上游任务
    \item \textbf{NetCDF 中间数据}：阶段产物以 \texttt{EnsembleData} NetCDF 文件存储，支持自描述、可追踪的中间结果
    \item \textbf{LLM 驱动执行}：每个阶段的任务由 LLM 代理自主调用工具完成
\end{itemize}

\subsection{五阶段分析工作流}

\begin{center}
\begin{tcolorbox}[width=0.9\textwidth,boxrule=0.5pt,colback=gray!5,colframe=gray!50]
\begin{verbatim}
1. correlator_analysis  →  stages/correlator/    （关联函数分析）
2. renormalization      →  stages/renorm/        （重整化）
3. fourier_transform    →  stages/fourier/       （傅里叶变换）
4. perturbative_matching →  stages/matching/      （微扰匹配）
5. extrapolation        →  stages/extrapolation/  （外推）
\end{verbatim}
\end{tcolorbox}
\end{center}

重整化阶段支持三种方案：
\begin{itemize}[itemsep=2pt]
    \item \textbf{比值方案（ratio）}：目标与分母在完整坐标网格上逐点相除
    \item \textbf{混合比值方案（hybrid-ratio）}：短距离比值 + 长距离指数修正
    \item \textbf{混合自重整化方案（hybrid-self-renormalization）}：全范围自重整化 + 短距离 $\overline{\text{MS}}$ 匹配
\end{itemize}

% ============================================================================
\section{安装指南}
% ============================================================================

\subsection{环境要求}

\begin{table}[h]
\centering
\begin{tabular}{@{}ll@{}}
\toprule
\textbf{组件} & \textbf{要求} \\
\midrule
Python       & $\ge$ 3.10 \\
\multicolumn{2}{@{}l}{\textbf{核心依赖}} \\
numpy        & $\ge$ 2.0 \\
pydantic     & $\ge$ 2.8 \\
typer        & $\ge$ 0.12 \\
tqdm         & $\ge$ 4.68 \\
\multicolumn{2}{@{}l}{\textbf{分析依赖（analysis extra）}} \\
gvar         & $\ge$ 12.1 \\
lsqfit       & $\ge$ 13.0 \\
scipy        & $\ge$ 1.13 \\
h5py         & $\ge$ 3.11 \\
matplotlib   & $\ge$ 3.8 \\
xarray       & $\ge$ 2024.1 \\
netCDF4      & $\ge$ 1.7 \\
\bottomrule
\end{tabular}
\caption{系统需求}
\end{table}

\subsection{安装步骤}

\begin{lstlisting}[style=bashcode]
# 1. 克隆仓库
git clone https://github.com/Greyyy-HJC/lamet-agent.git
cd lamet-agent

# 2. 创建虚拟环境
python3 -m venv .venv
source .venv/bin/activate

# 3. 安装核心包 + 开发和分析依赖
pip install -e ".[dev,analysis]"

# 4. （可选）安装 Codex 后端支持
pip install -e ".[codex]"
\end{lstlisting}

% ============================================================================
\section{清单（Manifest）配置}
% ============================================================================

清单是 lamet-agent 的核心配置机制，定义了分析运行的所有参数。

\subsection{顶层结构}

\begin{lstlisting}[style=jsoncode]
{
  "metadata": { ... },   // 运行级设置
  "inputs":   { ... },   // 数据源（关联函数、外部产物、核函数）
  "stages":   { ... }    // 阶段配置（默认值 + 任务列表）
}
\end{lstlisting}

\subsection{Metadata 字段}

\begin{table}[h]
\centering
\begin{tabular}{@{}>{\ttfamily}lll@{}}
\toprule
\textbf{字段} & \textbf{类型} & \textbf{说明} \\
\midrule
run\_id              & string  & 运行标识符 \\
root\_directory      & string  & 根目录（所有相对路径的基准）\\
artifacts\_directory & string  & 产物输出目录 \\
target\_observable   & string  & 目标可观测量："pdf" 或 "da" \\
resample\_mode       & string  & 重采样模式："jk", "bs", 或 "gvar" \\
random\_seed         & int     & 随机种子（必需）\\
sample\_error\_mode  & string  & 误差模式（默认 "covariance"）\\
bin\_size            & int     & 配置合并大小（可选）\\
workers              & int     & 并行工作进程数（默认 1）\\
stages               & list    & 有序的阶段 ID 列表 \\
\bottomrule
\end{tabular}
\caption{Metadata 参数}
\end{table}

\subsection{Inputs 配置}

\begin{lstlisting}[style=jsoncode]
{
  "inputs": {
    "correlators": [
      {
        "correlator_id": "ensemble_2pt",
        "correlator_type": "2pt",
        "data_path": "data/ensemble_2pt.h5",
        "ensemble": "HISQa060_X",
        "hadron": "pion",
        "gfix": "CG",
        "source_operator": "g5",
        "sink_operator": "g5",
        "volume": "S48T64",
        "lattice_spacing_fm": 0.0574,
        "momentum": ["PX0PY0PZ0", "PX5PY0PZ0"]
      }
    ],
    "artifacts": [
      {
        "id": "bare_pdf_reference",
        "stage": "correlator_analysis",
        "path": "artifacts/correlator_analysis/ref.nc"
      }
    ],
    "kernels": [
      {
        "stage": "renormalization",
        "kernel_id": "ZMSbar_pdf",
        "kernel_path": "src/lamet_agent/kernels.py",
        "scheme": "hybrid_ratio",
        "kernel_parameters": { "mu": 2.0 }
      }
    ]
  }
}
\end{lstlisting}

\subsection{动量标签格式}

动量标签采用精确格式 \texttt{PXnPYnPZn}（每个分量为有符号整数），物理动量由下式计算：

\[
p\,[\mathrm{GeV}] = \frac{2\pi\hbar c}{L_s\,a\,[\mathrm{fm}]} \sqrt{n_x^2+n_y^2+n_z^2},
\qquad \hbar c = 0.1973269804\ \mathrm{GeV\,fm}
\]

\subsection{标准 HDF5 关联函数格式}

\begin{itemize}[itemsep=3pt]
    \item \textbf{普通 2pt}：\texttt{<source>/<sink>/<momentum>}，形状 \texttt{(Lt, n\_cfg)}
    \item \textbf{qDA 2pt}：\texttt{<source>/<sink>/<momentum>/bT<bT>/bz<bz>}，形状 \texttt{(Lt, n\_cfg)}
    \item \textbf{3pt}：\texttt{<source>/<sink>/<current>/<momentum>/tsep<tsep>/bT<bT>/bz<bz>}，形状 \texttt{(tsep+1, n\_cfg)}
\end{itemize}

% ============================================================================
\section{命令行接口}
% ============================================================================

\subsection{CLI 命令}

\texttt{lamet-agent} 通过 \texttt{pyproject.toml} 注册为命令行工具：

\begin{lstlisting}[style=bashcode]
# 验证清单
lamet-agent validate examples/pion_pdf_cg_manifest.json

# 运行分析
lamet-agent run examples/pion_pdf_cg_manifest.json

# 交互式规划清单
lamet-agent plan draft_manifest.jsonc --backend api --model deepseek/deepseek-chat
\end{lstlisting}

\subsection{run 命令参数}

\begin{table}[h]
\centering
\begin{tabular}{@{}>{\ttfamily}ll@{}}
\toprule
\textbf{参数} & \textbf{说明} \\
\midrule
--backend & LLM 后端: mock, external, api, codex \\
--model   & 模型 ID（api: provider/model\_id, codex: CODEX\_MODEL\_ID）\\
--verbose / -v & 启用 ReAct 风格详细输出 \\
--actions-path & external 后端的 JSONL 转录路径 \\
--api-key-file & API 密钥文件路径（默认 api.key）\\
--base-url & 覆盖 API HTTP 端点 \\
--report\_language & 报告语言: en 或 ch \\
\bottomrule
\end{tabular}
\caption{run 命令参数}
\end{table}

\subsection{运行示例}

\begin{lstlisting}[style=bashcode]
# 使用 DeepSeek API 后端运行
lamet-agent run examples/pion_pdf_cg_manifest.json \
    --backend api --model deepseek/deepseek-chat --verbose

# 使用 OpenAI 后端运行
lamet-agent run examples/pion_pdf_cg_manifest.json \
    --backend api --model openai/gpt-4o

# 使用 Codex 后端运行
lamet-agent run examples/pion_pdf_cg_manifest.json \
    --backend codex --model CODEX_MODEL_ID --verbose

# Mock 模式（开发/测试）
lamet-agent run examples/pion_pdf_cg_manifest.json --backend mock

# 重放 JSONL 转录
lamet-agent run examples/pion_pdf_cg_manifest.json \
    --backend external --actions-path actions.jsonl
\end{lstlisting}

% ============================================================================
\section{项目结构}
% ============================================================================

\begin{lstlisting}[style=bashcode]
lamet-agent/
├── pyproject.toml                  # 项目配置、依赖、入口点
├── README.md                       # 完整文档
├── PLAN.md / AGENTS.md / PROJECT_LOG.md  # 开发规划与日志
├── examples/                       # 示例清单
│   ├── sample_manifest.jsonc       # 带注释的参考清单（JSONC）
│   ├── pion_pdf_cg_manifest.json   # P0/P5 CG pion PDF
│   ├── pion_pdf_gi_manifest.json   # P0/P4 GI pion PDF
│   ├── pion_da_gi_manifest.json    # GI pion DA
│   ├── kaon_da_gi_manifest.json    # GI kaon DA
│   ├── partial_sample_manifest.jsonc
│   └── fake_data/generate_fake_data.py
├── src/lamet_agent/                # 核心包
│   ├── __init__.py
│   ├── agent.py                    # run_agent() 主循环
│   ├── cli.py                      # CLI 入口（validate, run, plan）
│   ├── manifest.py                 # 清单模式与验证
│   ├── manifest_params.py          # 阶段参数合约注册表
│   ├── kernels.py                  # 内置核函数示例
│   ├── core/                       # 核心模块
│   │   ├── data.py                 # EnsembleData/EnsembleInfo 类型
│   │   ├── llm.py                  # 可插拔 LLM 会话后端
│   │   ├── prompting.py            # 提示词构建
│   │   ├── tools.py                # 工具注册、参数准备、绘图路径
│   │   ├── stages.py               # 阶段 ID → 包映射
│   │   ├── trace.py                # ReAct 风格追踪
│   │   ├── banner.py               # 启动横幅
│   │   ├── plotting.py             # 出版风格绑图
│   │   ├── resampling.py           # Jackknife/Bootstrap 重采样
│   │   └── reporting.py            # 阶段报告
│   ├── planning/                   # 清单规划模式
│   │   ├── agent.py                # 规划代理
│   │   ├── core.py                 # 规划核心逻辑
│   │   ├── conversion.py           # HDF5 数据转换
│   │   ├── questions.py            # 交互式问题
│   │   └── render.py               # 终端渲染
│   └── stages/                     # 阶段实现
│       ├── correlator/             # 关联函数分析
│       ├── renorm/                 # 重整化
│       ├── fourier/                # 傅里叶变换
│       ├── matching/               # 微扰匹配
│       ├── extrapolation/          # 外推
│       └── review/                 # 审查
│           (每个阶段含 prompts.py, skills.py, functions.py, reporting.py)
└── tests/unit/                     # 单元测试
\end{lstlisting}

% ============================================================================
\section{核心模块详解}
% ============================================================================

\subsection{代理工作流（agent.py）}

\begin{lstlisting}[style=pycode]
def run_agent(manifest_path, *, backend, model=None,
              verbose=False, actions_path=None,
              api_key_file=None, base_url=None,
              report_language="en") -> dict:
    """
    执行完整的清单驱动分析工作流。

    返回: {"run_id": ..., "status": ..., "summary": ...,
            "stage_results": ..., "actions": ...}
    """
\end{lstlisting}

\textbf{执行流程}：
\begin{enumerate}[itemsep=3pt]
    \item CLI 接收清单路径和运行时选项
    \item \texttt{manifest.py} 验证源 ID、任务 ID、有序依赖和路径
    \item \texttt{agent.py} 按序执行 \texttt{metadata.stages} 列表
    \item 对每个阶段任务：
    \begin{enumerate}
        \item \texttt{core/tools.validate\_stage\_inputs()} 检查缺失输入
        \item \texttt{core/prompting.build\_stage\_static\_prompt()} 组装静态上下文
        \item \texttt{core/llm.make\_llm\_session()} 创建 LLM 会话
        \item 多轮循环（最多 40 个工具步骤）：模型发出 JSON 动作，工具执行并返回观察
        \item 终端工具将主要数据放入 \texttt{store["output"]}
    \end{enumerate}
    \item 阶段完成后生成报告
    \item 输出紧凑的 JSON 摘要到 stdout
\end{enumerate}

\subsection{LLM 后端（core/llm.py）}

支持四种可插拔会话后端：

\begin{table}[h]
\centering
\begin{tabular}{@{}lll@{}}
\toprule
\textbf{后端} & \textbf{说明} & \textbf{用途} \\
\midrule
\texttt{mock}     & 确定性脚手架              & 开发/测试 \\
\texttt{external} & JSONL 转录重放            & 回归测试 \\
\texttt{api}      & OpenAI 兼容 HTTP API      & 生产运行 \\
\texttt{codex}    & Codex Python SDK          & Codex 集成 \\
\bottomrule
\end{tabular}
\caption{LLM 后端类型}
\end{table}

\begin{lstlisting}[style=pycode]
class LlmSession:
    """可插拔 LLM 会话抽象"""

def make_llm_session(backend, *, model=None, api_key_file=None,
                     base_url=None, actions_path=None) -> LlmSession:
    """工厂函数：根据后端类型创建会话"""

def parse_api_model(spec: str) -> tuple[str, str]:
    """解析 provider/model_id CLI 规格"""

# 预配置的 API 提供商
PROVIDERS = {
    "deepseek": {"base_url": "https://api.deepseek.com/v1",
                  "default_model": "deepseek-chat",
                  "api_key_env": "DEEPSEEK_API_KEY"},
    "openai":   {"base_url": "https://api.openai.com/v1",
                  "default_model": "gpt-4o",
                  "api_key_env": "OPENAI_API_KEY"},
}
\end{lstlisting}

\subsection{数据类型（core/data.py）}

\begin{lstlisting}[style=pycode]
class EnsembleInfo:
    """系综元数据容器"""
    ensemble: str
    hadron: str
    gfix: str
    momentum: str
    volume: str
    lattice_spacing_fm: float
    # ... 更多字段

class EnsembleData:
    """重采样格点数据容器（NetCDF 序列化）"""
    values: np.ndarray       # shape (n_sample, *physical_dims)
    coords: dict             # 物理坐标
    ensemble: EnsembleInfo   # 系综信息
    resample: str            # 重采样模式

    def to_netcdf(self, path: str) -> None
    @classmethod
    def from_netcdf(cls, path: str) -> EnsembleData
\end{lstlisting}

\textbf{NetCDF 数据结构}：
\begin{itemize}[itemsep=2pt]
    \item 第一维度始终命名为 \texttt{resample}
    \item 复数数组原生往返（\texttt{auto\_complex=True}）
    \item 属性：\texttt{ensemble}（JSON 序列化的 EnsembleInfo）、\texttt{resample}
\end{itemize}

\subsection{工具系统（core/tools.py）}

\begin{lstlisting}[style=pycode]
def prepare_tool_args(stage_name, tool_name, tool_args, manifest) -> dict:
    """规范化 LLM 工具调用参数（路径、保存路径等）"""

def filter_tool_kwargs(tool_func, tool_args, *, strict=True) -> dict:
    """过滤工具参数以匹配函数签名"""

def resolve_stage_tools(stage_name, job_store) -> dict:
    """解析阶段的工具注册表"""

def resolve_plot_save_path(manifest, stage_name, job_id, stem) -> str:
    """解析绘图保存路径到阶段产物目录"""

def validate_stage_inputs(stage_name, job, job_store,
                          manifest_inputs) -> dict:
    """验证阶段输入并返回 input_issues"""
\end{lstlisting}

\subsection{清单系统（manifest.py）}

\begin{lstlisting}[style=pycode]
def load_manifest(path: str) -> dict:
    """加载和解析清单文件（JSON 或 JSONC）"""

def validate_manifest(manifest: dict) -> list[str]:
    """验证清单模式、任务 ID、有序依赖和路径"""

def resolve_manifest_paths(manifest: dict, manifest_dir: str) -> dict:
    """解析所有相对路径"""
\end{lstlisting}

\subsection{阶段参数合约（manifest\_params.py）}

\begin{lstlisting}[style=pycode]
# 中央参数合约注册表
STAGE_PARAM_CONTRACS = {
    "correlator_analysis": { ... },
    "renormalization":     { ... },
    "fourier_transform":   { ... },
    "perturbative_matching": { ... },
    "extrapolation":       { ... },
}

def validate_stage_params(stage_name, defaults, job_params) -> list[str]:
    """递归拒绝未知键（defaults 和 params）"""
\end{lstlisting}

% ============================================================================
\section{各阶段详解}
% ============================================================================

\subsection{关联函数分析（correlator\_analysis）}

\textbf{工具集}：
\begin{itemize}[itemsep=3pt]
    \item \texttt{inspect\_correlator\_scale}：选择关联函数缩放因子
    \item \texttt{tune\_ground\_state}：2pt 窗口扫描 + 模型平均
    \item \texttt{tune\_bare\_matrix}：在样本平均数据上扫描裸矩阵元拟合窗口
    \item \texttt{fit\_bare\_matrix\_grid}：将调优后的窗口应用于所有 $z$ 和重采样样本
\end{itemize}

\textbf{参数}：
\begin{itemize}[itemsep=2pt]
    \item \texttt{fit\_scope}：分析族（\texttt{3pt\_ratio}, \texttt{FH}, \texttt{3pt\_ratio+FH}, \texttt{qda\_ratio}）
    \item \texttt{model\_average}：是否对拟合函数候选进行 BMA 组合
    \item \texttt{fit\_strategy}：\texttt{joint}, \texttt{chained}, 或 \texttt{independent}
\end{itemize}

\subsection{重整化（renormalization）}

\textbf{方案}：
\begin{itemize}[itemsep=3pt]
    \item \texttt{ratio}：$h_s^R(z) = h_s^{\text{target}}(z) / h_s^{\text{denominator}}(z)$
    \item \texttt{hybrid\_ratio}：短距离比值 + 长距离指数修正
    \item \texttt{hybrid\_self\_renormalization}：全范围自重整化 + $\overline{\text{MS}}$ 匹配
\end{itemize}

\textbf{混合自重整化工作流}：
\[
g(z)-\ln Z_{\overline{\mathrm{MS}}}^{\mathrm{PDF}}(z;\mu)\simeq m_0z+b,
\qquad
z_R(z,a)=\exp[\ln M_{\mathrm{fit}}(z,a)-g(z)+m_0z]
\]
\[
H_{\mathrm{ren}}(z) = \frac{H_{\mathrm{bare}}(z)}{z_R(z,a)\,Z_{\overline{\mathrm{MS}}}(z)}
\]

\textbf{参数}：
\begin{itemize}[itemsep=2pt]
    \item \texttt{normalization}：是否在重整化前除以 $z=0$ 值
    \item \texttt{zs\_fm}：混合开关距离（fm）
    \item \texttt{scheme\_parameters.LambdaQCD\_gev}：$\Lambda_{\mathrm{QCD}}$（GeV）
    \item \texttt{scheme\_parameters.d}：连续/离散化系数
    \item \texttt{scheme\_parameters.m0\_gev}：$m_0$ 参数
    \item \texttt{mu}：重整化标度（GeV）
\end{itemize}

\subsection{傅里叶变换（fourier\_transform）}

将坐标空间的矩阵元变换为准 PDF：
\[
\tilde{g}(x, P_z) = \int \frac{dz}{2\pi} e^{ix z P_z} h^R(z, P_z)
\]

\subsection{微扰匹配（perturbative\_matching）}

将准 PDF 匹配到光锥 PDF：
\[
g(x, \mu) = \int \frac{dy}{|y|} C\left(\frac{x}{y}, \frac{\mu}{P_z}\right) \tilde{g}(y, P_z)
\]

\subsection{外推（extrapolation）}

将有限动量下的结果外推到无限动量极限。

% ============================================================================
\section{中间数据格式（NetCDF）}
% ============================================================================

\subsection{产物链}

\begin{table}[h]
\centering
\begin{tabular}{@{}ll@{}}
\toprule
\textbf{阶段} & \textbf{示例产物} \\
\midrule
\texttt{correlator\_analysis} & \texttt{correlator\_analysis/ca\_p5.nc} \\
\texttt{renormalization}     & \texttt{renormalization/rn\_p5.nc} \\
\texttt{fourier\_transform}  & \texttt{fourier\_results/fourier\_result.nc} \\
\texttt{perturbative\_matching} & \texttt{matching\_results/quasi\_pdf.nc} \\
\bottomrule
\end{tabular}
\caption{典型产物链}
\end{table}

\subsection{Python 读写}

\begin{lstlisting}[style=pycode]
from lamet_agent.core.data import EnsembleData

# 写入
data.to_netcdf("artifacts/fourier_results/fourier_result.nc")

# 读取
reload = EnsembleData.from_netcdf("artifacts/fourier_results/fourier_result.nc")
\end{lstlisting}

\subsection{外部工具读取}

\begin{lstlisting}[style=bashcode]
# 使用 ncdump 查看
ncdump -h artifacts/fourier_results/fourier_result.nc
\end{lstlisting}

\begin{lstlisting}[style=pycode]
# 使用 xarray 读取（无需 lamet-agent）
import json, xarray as xr
from lamet_agent.core.data import EnsembleInfo

da = xr.load_dataarray("fourier_result.nc", auto_complex=True)
ensemble = EnsembleInfo(**json.loads(da.attrs["ensemble"]))
values = da.values  # shape (n_sample, *physical_dims)
\end{lstlisting}

% ============================================================================
\section{示例清单详解}
% ============================================================================

\subsection{pion\_pdf\_cg\_manifest.json —— Pion PDF（Coulomb 规范）}

完整的 P0/P5 工作流，包含关联函数分析、混合比值重整化、傅里叶变换和微扰匹配。

\subsection{pion\_da\_gi\_manifest.json —— Pion DA（规范不变）}

完整的 GI pion DA 工作流，从 qDA 关联函数分析到匹配和审查。

\subsection{kaon\_da\_gi\_manifest.json —— Kaon DA（规范不变）}

完整的 GI kaon DA 工作流，运行辅助脚本位于 \texttt{runs/ds\_kaon\_da\_gi/}。

\subsection{sample\_manifest.jsonc —— 带注释参考清单}

标记为 JSONC（JSON with comments），每个字段都带有文档注释。

% ============================================================================
\section{规划模式（Plan Mode）}
% ============================================================================

\texttt{lamet-agent plan} 命令接受不完整的 JSON/JSONC 清单，运行 LLM 控制的规划循环：

\begin{lstlisting}[style=bashcode]
lamet-agent plan draft_manifest.jsonc --backend api --model deepseek/deepseek-chat
\end{lstlisting}

\textbf{规划循环功能}：
\begin{itemize}[itemsep=2pt]
    \item 检查清单完整性
    \item 检查 HDF5 输入
    \item 提出终端问题
    \item 应用经验证的 JSON Patch 编辑
    \item 生成快速清单（jackknife + 均值误差）和完整清单（协方差 + BMA）
    \item 自动转换非标准 HDF5 布局
\end{itemize}

\textbf{终端摘要组织为}：缺失参数、不一致设置、建议修改、数据转换。

% ============================================================================
\section{接口清单（API Reference）}
% ============================================================================

\subsection{CLI 入口（cli.py）}

\begin{lstlisting}[style=pycode]
def entrypoint():
    """CLI 入口点：注册 validate, run, plan 命令"""

def cmd_validate(manifest_path: str):
    """验证清单模式与核函数导入"""

def cmd_run(manifest_path, *, backend, model, verbose,
            actions_path, api_key_file, base_url, report_language):
    """运行完整的阶段循环并收集结构化动作"""

def cmd_plan(draft_path, *, backend, model, api_key_file, base_url):
    """交互式规划清单"""
\end{lstlisting}

\subsection{代理核心（agent.py）}

\begin{lstlisting}[style=pycode]
def run_agent(manifest_path, *, backend, model=None,
              verbose=False, actions_path=None,
              api_key_file=None, base_url=None,
              report_language="en") -> dict:
    """
    执行 metadata.stages，按序运行每个声明的任务，
    使用隔离的 store，将 store["output"] 注册到任务 ID 下。
    """
\end{lstlisting}

\subsection{数据核心（core/data.py）}

\begin{lstlisting}[style=pycode]
@dataclass
class EnsembleInfo:
    """系综元数据"""
    ensemble: str
    hadron: str = ""
    gfix: str = ""
    momentum: str = ""
    volume: str = ""
    lattice_spacing_fm: float = 0.0
    bz_direction: str = ""
    # ... 更多字段

@dataclass
class EnsembleData:
    """重采样格点数据"""
    values: np.ndarray
    coords: dict[str, np.ndarray]
    ensemble: EnsembleInfo
    resample: str = "jk"
    attrs: dict = field(default_factory=dict)

    def to_netcdf(self, path: str) -> None: ...
    @classmethod
    def from_netcdf(cls, path: str) -> EnsembleData: ...
\end{lstlisting}

\subsection{LLM 会话（core/llm.py）}

\begin{lstlisting}[style=pycode]
class LlmSession(ABC):
    @abstractmethod
    def send(self, messages: list[dict]) -> dict: ...

class MockSession(LlmSession): ...
class ExternalSession(LlmSession): ...
class ApiSession(LlmSession): ...
class CodexSession(LlmSession): ...

def make_llm_session(backend, **kwargs) -> LlmSession: ...
def parse_api_model(spec: str) -> tuple[str, str]: ...
\end{lstlisting}

\subsection{阶段映射（core/stages.py）}

\begin{lstlisting}[style=pycode]
STAGE_PACKAGES = {
    "correlator_analysis":     "lamet_agent.stages.correlator",
    "renormalization":         "lamet_agent.stages.renorm",
    "fourier_transform":       "lamet_agent.stages.fourier",
    "perturbative_matching":   "lamet_agent.stages.matching",
    "extrapolation":           "lamet_agent.stages.extrapolation",
}
\end{lstlisting}

\subsection{关联函数阶段工具（stages/correlator/functions.py）}

\begin{lstlisting}[style=pycode]
def inspect_correlator_scale(...) -> dict:
    """选择关联函数缩放因子"""

def tune_ground_state(...) -> dict:
    """2pt 窗口扫描 + 模型平均"""

def tune_bare_matrix(...) -> dict:
    """在样本平均数据上扫描裸矩阵元拟合窗口"""

def fit_bare_matrix_grid(...) -> dict:
    """将统一调优的窗口应用于所有 z 和重采样样本"""

STAGE_TOOLS = {
    "inspect_correlator_scale": inspect_correlator_scale,
    "tune_ground_state": tune_ground_state,
    "tune_bare_matrix": tune_bare_matrix,
    "fit_bare_matrix_grid": fit_bare_matrix_grid,
}
\end{lstlisting}

\subsection{重整化阶段工具（stages/renorm/functions.py）}

\begin{lstlisting}[style=pycode]
def apply_ratio_scheme_renormalization(...) -> dict:
    """应用比值方案重整化"""

def apply_hybrid_ratio_renormalization(...) -> dict:
    """应用混合比值方案"""

def fit_self_renormalization_factor(...) -> dict:
    """拟合自重整化因子"""

def apply_self_renormalization(...) -> dict:
    """应用自重整化因子"""

def plot_self_renormalization_diagnostics(...) -> dict:
    """绑制自重整化诊断图"""

def plot_renormalized_matrix_element(...) -> dict:
    """绑制重整化矩阵元图"""
\end{lstlisting}

\subsection{傅里叶阶段工具（stages/fourier/functions.py）}

\begin{lstlisting}[style=pycode]
def fourier_transform(...) -> dict:
    """执行傅里叶变换（坐标空间 → 动量空间）"""

def plot_fourier_result(...) -> dict:
    """绑制傅里叶变换结果"""
\end{lstlisting}

\subsection{匹配阶段工具（stages/matching/functions.py）}

\begin{lstlisting}[style=pycode]
def apply_matching(...) -> dict:
    """应用微扰匹配核函数"""

def plot_matching_result(...) -> dict:
    """绑制匹配结果"""
\end{lstlisting}

\subsection{外推阶段工具（stages/extrapolation/functions.py）}

\begin{lstlisting}[style=pycode]
def extrapolate(...) -> dict:
    """执行无限动量外推"""

def plot_extrapolation(...) -> dict:
    """绑制外推结果"""
\end{lstlisting}

\subsection{核函数（kernels.py）}

\begin{lstlisting}[style=pycode]
# 内置核函数示例，用于烟雾测试
# 生产核函数通过清单中的 kernel_path 指定
\end{lstlisting}

% ============================================================================
\section{依赖关系}
% ============================================================================

\begin{table}[h]
\centering
\begin{tabular}{@{}lll@{}}
\toprule
\textbf{依赖}    & \textbf{版本}    & \textbf{用途} \\
\midrule
Python           & $\ge$ 3.10       & 运行环境 \\
numpy            & $\ge$ 2.0        & 数值计算 \\
pydantic         & $\ge$ 2.8        & 数据验证与模式 \\
typer            & $\ge$ 0.12       & CLI 框架 \\
tqdm             & $\ge$ 4.68       & 进度条 \\
gvar             & $\ge$ 12.1       & 误差传播（分析 extra）\\
lsqfit           & $\ge$ 13.0       & 最小二乘拟合（分析 extra）\\
scipy            & $\ge$ 1.13       & 科学计算（分析 extra）\\
h5py             & $\ge$ 3.11       & HDF5 I/O（分析 extra）\\
matplotlib       & $\ge$ 3.8        & 绑图（分析 extra）\\
xarray           & $\ge$ 2024.1     & 标注数组（分析 extra）\\
netCDF4          & $\ge$ 1.7        & NetCDF I/O（分析 extra）\\
openai-codex     & $\ge$ 0.1.0b3    & Codex 后端（codex extra）\\
\bottomrule
\end{tabular}
\caption{完整依赖列表}
\end{table}

% ============================================================================
\section{报告系统}
% ============================================================================

每个阶段完成后，\texttt{reporting.py} 模块自动生成指定语言（en/ch）的分析报告，包含：
\begin{itemize}[itemsep=2pt]
    \item 阶段执行摘要
    \item 工具调用日志
    \item 拟合参数与结果
    \item 绑图输出路径
    \item 中间数据产物路径
\end{itemize}

\end{document}
